\notechapter{N-20260417}{Green's Theorem, Conservative Fields, and the Planar Stokes Principle}

\section{Green's theorem}

Green's theorem relates circulation around a boundary to local rotation inside the region.

\begin{theorem}[Green's theorem]
Let \(D\subseteq\mathbb R^2\) be a positively oriented region with piecewise smooth simple closed boundary \(C=\partial D\).  If \(P\) and \(Q\) have continuous first partial derivatives on a neighborhood of \(D\), then
\[
  \oint_C P\,dx+Q\,dy
  =\iint_D \left(\frac{\partial Q}{\partial x}-\frac{\partial P}{\partial y}\right)\,dA.
\]
\end{theorem}

The left side is circulation along the boundary.  The right side is total scalar curl over the region.

\begin{figure}[H]
\centering
\includegraphics[width=0.82\textwidth,height=0.74\textheight,keepaspectratio]{figures/notes-md/N-20260417/N-20260417-fig-01-green-simple-closed-curves.png}
\caption{Simple closed curves and non-simple closed curves in the statement of Green's theorem.}
\end{figure}

\section{Proof idea: rectangles and cancellation}

For a rectangle \([a,b]\times[c,d]\), the theorem follows from the fundamental theorem of calculus.  The \(Q\,dy\) part detects horizontal change in \(Q\), and the \(P\,dx\) part detects vertical change in \(P\) with the opposite sign.

For a region subdivided into many rectangles, internal edges are traversed twice in opposite directions.  Their line integrals cancel.  Only the outer boundary remains.

\begin{figure}[H]
\centering
\includegraphics[width=0.82\textwidth,height=0.74\textheight,keepaspectratio]{figures/notes-md/N-20260417/N-20260417-fig-03-green-vertical-simple-region-proof.png}
\caption{For a vertically simple region, Green's theorem is proved by reducing to one-variable fundamental theorem computations.}
\end{figure}

\begin{figure}[H]
\centering
\includegraphics[width=0.82\textwidth,height=0.74\textheight,keepaspectratio]{figures/notes-md/N-20260417/N-20260417-fig-04-green-subdivision-boundary-cancellation.png}
\caption{Subdivision by auxiliary arcs: internal boundary contributions cancel, leaving only the outer boundary.}
\end{figure}

\section{Conservative fields and path independence}

A vector field \(F=(P,Q)\) is \vocab{conservative} on a region \(U\) if there is a potential function \(\phi\) such that
\[
  \nabla\phi=F.
\]
Then
\[
  \int_C P\,dx+Q\,dy=\phi(B)-\phi(A)
\]
for every path \(C\) from \(A\) to \(B\).  Thus the integral depends only on endpoints.

\begin{proposition}
If \(F=\nabla\phi\), then \(P_y=Q_x\).
\end{proposition}

\begin{proof}
If \(P=\phi_x\) and \(Q=\phi_y\), then under the usual continuity hypotheses,
\[
  P_y=\phi_{xy}=\phi_{yx}=Q_x.
\]
\end{proof}

The converse requires a topological hypothesis.

\begin{theorem}[Conservative-field test]
If \(U\subseteq\mathbb R^2\) is simply connected and \(P_y=Q_x\) on \(U\), then \(F=(P,Q)\) is conservative on \(U\).
\end{theorem}

\begin{figure}[H]
\centering
\includegraphics[width=0.82\textwidth,height=0.74\textheight,keepaspectratio]{figures/notes-md/N-20260417/N-20260417-fig-02-green-simply-connected-regions.png}
\caption{Simply connected and multiply connected plane regions.  The distinction matters for path independence and conservative fields.}
\end{figure}

\section{The punctured-plane warning}

The standard example is
\[
  F(x,y)=\left(-\frac{y}{x^2+y^2},\frac{x}{x^2+y^2}\right)
\]
on \(\mathbb R^2\setminus\{0\}\).  A direct computation gives \(P_y=Q_x\) away from the origin, but around the unit circle \(r(t)=(\cos t,\sin t)\),
\[
  F(r(t))=(-\sin t,\cos t),\qquad r'(t)=(-\sin t,\cos t),
\]
so
\[
  \oint_{S^1}F\cdot dr=\int_0^{2\pi}1\,dt=2\pi.
\]
Therefore no global potential exists on the punctured plane.  The hole is detected by circulation.

\section{Area by Green's theorem}

Green's theorem can compute area.  Choose
\[
  P=-\frac{y}{2},\qquad Q=\frac{x}{2}.
\]
Then \(Q_x-P_y=1\), so
\[
  \operatorname{area}(D)=\iint_D1\,dA
  =\frac{1}{2}\oint_{\partial D}x\,dy-y\,dx.
\]

\begin{example}
For the unit circle \(x=\cos t\), \(y=\sin t\),
\[
  \frac{1}{2}\int_0^{2\pi}(x y'-y x')\,dt
  =\frac{1}{2}\int_0^{2\pi}1\,dt=\pi.
\]
\end{example}

\section{Differential forms notation}

Let
\[
  \omega=P\,dx+Q\,dy.
\]
Then
\[
  d\omega=(Q_x-P_y)\,dx\wedge dy.
\]
Green's theorem becomes
\[
  \int_{\partial D}\omega=\int_D d\omega.
\]
This is Stokes' theorem in the plane.

The algebraic slogan is: integrate a form over the boundary, or integrate its derivative over the interior.

\section{Closed forms, exact forms, and topology}

A one-form \(\omega\) is \vocab{closed} if \(d\omega=0\).  It is \vocab{exact} if \(\omega=d\phi\) for some function \(\phi\).  Exact forms are always closed, because \(d^2=0\).  The converse depends on the domain.

On simply connected plane regions, closed one-forms are exact.  On the punctured plane, the form
\[
  \omega=\frac{-y\,dx+x\,dy}{x^2+y^2}
\]
is closed but not exact, because
\[
  \int_{S^1}\omega=2\pi.
\]

\section{Connection with de Rham cohomology}

The first de Rham cohomology group is
\[
  H^1_{\mathrm{dR}}(M)
  =\frac{\ker(d:\Omega^1(M)\to\Omega^2(M))}
    {\operatorname{im}(d:\Omega^0(M)\to\Omega^1(M))}.
\]
It measures closed one-forms modulo exact one-forms.  In elementary vector calculus language, it measures the obstruction to finding potentials for curl-free fields.

\section{Winding number hidden in the punctured-plane example}

For the one-form
\[
  \omega=\frac{-y\,dx+x\,dy}{x^2+y^2},
\]
write \(x=r\cos\theta\), \(y=r\sin\theta\).  Then formally
\[
  \omega=d\theta.
\]
The word "formally" matters: \(\theta\) is not a globally single-valued function on \(\mathbb R^2\setminus\{0\}\).  Around a closed loop \(\gamma\),
\[
  \frac{1}{2\pi}\int_\gamma \omega
\]
is the winding number of \(\gamma\) about the origin.

For the unit circle, the winding number is \(1\).  For a loop going twice around, it is \(2\).  For a loop not enclosing the origin, it is \(0\).  Thus the line integral is not merely detecting that a potential is missing; it measures how the curve winds around the hole.

\section{From Green's theorem to residues}

The same punctured-plane form is connected to complex analysis.  Since
\[
  \frac{dz}{z}=\frac{dx+i\,dy}{x+iy},
\]
one computes
\[
  \operatorname{Im}\frac{dz}{z}
  =\frac{x\,dy-y\,dx}{x^2+y^2}=\omega.
\]
Thus
\[
  \int_{S^1}\omega=2\pi
\]
is the imaginary part of
\[
  \int_{S^1}\frac{dz}{z}=2\pi i.
\]
The residue theorem is a complex-analytic refinement of the same topological fact: a closed integral can detect a puncture and its winding number.

\section{Cohomology as bookkeeping for failed potentials}

The statement "curl-free implies conservative" is not universally true; it is true when the relevant cohomology vanishes.  In de Rham notation,
\[
  H^1_{\mathrm{dR}}(M)
  =\frac{\text{closed one-forms}}{\text{exact one-forms}}.
\]
A nonzero class is represented by a closed form that is not the differential of any global function.

For \(M=\mathbb R^2\setminus\{0\}\), the class of \(\omega\) above generates \(H^1_{\mathrm{dR}}(M)\cong\mathbb R\).  The integral over \(S^1\) is a period pairing:
\[
  ([\omega],[S^1])\longmapsto \int_{S^1}\omega.
\]
So the elementary warning about holes is really the first example of pairing cohomology with cycles.

\section{Boundary cancellation and the algebra \texorpdfstring{$\partial^2=0$}{boundary squared is zero}}

The rectangle-subdivision proof of Green's theorem says that every internal edge appears twice with opposite orientations.  Algebraically this is the principle
\[
  \partial^2=0:
  \quad \text{the boundary of a boundary cancels.}
\]
De Rham theory has the dual identity
\[
  d^2=0.
\]
Stokes' theorem connects them:
\[
  \int_{\partial C}\omega=\int_C d\omega.
\]
If one applies this again, both sides become zero for structural reasons.  Thus Green's theorem is not an isolated vector-calculus trick; it is the two-dimensional visible face of a chain/cochain duality.

\section{What Green's theorem checks in practice}

Green's theorem converts a boundary integral into an area integral:
\[
  \oint_{\partial D} P\,dx+Q\,dy
  =\iint_D \left(\frac{\partial Q}{\partial x}-\frac{\partial P}{\partial y}\right)dA.
\]
In a computation this requires more than recognizing the formula.  The boundary must be positively oriented, the vector field must be defined on a region containing \(D\), and holes in the domain must be accounted for by extra boundary components with the correct orientation.

The punctured-plane example explains why this check matters.  The form
\[
  \omega=\frac{-y\,dx+x\,dy}{x^2+y^2}
\]
is closed away from the origin, but its integral around \(S^1\) is nonzero, so it is not globally exact on \(\mathbb R^2\setminus\{0\}\).  Green's theorem, the identity \(d^2=0\), and the cancellation of internal boundaries are therefore three views of the same mechanism: local derivative data controls boundary circulation only after the topology of the region has been included.
